\subsection{Diccionarios: dict, defaultdict y Counter}
\label{alg:2-3}

Los diccionarios son una de las estructuras de datos más utilizadas en programación competitiva. 
Permiten almacenar pares \texttt{clave : valor} con acceso promedio en tiempo constante.

En Python, la implementación utiliza una tabla hash, por lo que la mayoría de operaciones tienen complejidad promedio
\bigO{1} aunque en el peor caso pueden degradarse a \bigO{n}

si existen muchas colisiones.

\vspace{0.3cm}

\subsubsection{Crear un diccionario}

\begin{lstlisting}[language=Python]
freq = {}
persona = {"nombre": "David", "edad": 20}
\end{lstlisting}

\vspace{0.3cm}

\subsubsection{Acceso y recuperación}

Acceder directamente mediante corchetes genera un \texttt{KeyError} si la clave no existe.

\begin{lstlisting}[language=Python]
x = freq["a"]
\end{lstlisting}

Por ello normalmente es preferible utilizar \texttt{get()}.

\begin{lstlisting}[language=Python]
freq.get("a")
freq.get("a", 0)
freq.get("a", [])
\end{lstlisting}

\vspace{0.2cm}

Las funciones principales de recuperación son:

\begin{center}
\begin{tabular}{|l|p{6cm}|}
\hline
Función & Descripción \\ \hline
get() & Obtiene un valor sin lanzar excepción. \\ \hline
keys() & Devuelve todas las claves. \\ \hline
values() & Devuelve todos los valores. \\ \hline
items() & Devuelve pares (clave, valor). \\ \hline
\end{tabular}
\end{center}

Ejemplos:

\begin{lstlisting}[language=Python]
for k in freq.keys():
    ...
for v in freq.values():
    ...
for k, v in freq.items():
    ...
\end{lstlisting}

\vspace{0.3cm}

\subsubsection{Modificar un diccionario}

Insertar o modificar una clave.

\begin{lstlisting}[language=Python]
freq["a"] = 5
\end{lstlisting}

Actualizar múltiples valores.

\begin{lstlisting}[language=Python]
a = {"x":1}; b = {"y":2}
a.update(b)
\end{lstlisting}

Crear una clave únicamente si no existe.

\begin{lstlisting}[language=Python]
freq.setdefault("python", [])
freq["python"].append(10)
\end{lstlisting}

\vspace{0.3cm}

\subsubsection{Eliminar elementos}

Eliminar una clave.

\begin{lstlisting}[language=Python]
freq.pop("a")
\end{lstlisting}

Eliminar devolviendo un valor por defecto.

\begin{lstlisting}[language=Python]
freq.pop("a", 0)
\end{lstlisting}

Eliminar el último elemento insertado.

\begin{lstlisting}[language=Python]
freq.popitem()
\end{lstlisting}

Vaciar completamente el diccionario.

\begin{lstlisting}[language=Python]
freq.clear()
\end{lstlisting}

\vspace{0.3cm}

\subsubsection{Funciones integradas}

Cantidad de elementos.

\begin{lstlisting}[language=Python]
len(freq)
\end{lstlisting}

Copiar un diccionario.

\begin{lstlisting}[language=Python]
nuevo = freq.copy()
\end{lstlisting}

Comprobar existencia.

\begin{lstlisting}[language=Python]
if "python" in freq:
    ...
\end{lstlisting}

Eliminar una copia sin modificar el original.

\begin{lstlisting}[language=Python]
nuevo = dict(freq)
\end{lstlisting}

\vspace{0.3cm}

\subsubsection{Ordenar un diccionario}

Ordenar por clave.

\begin{lstlisting}[language=Python]
sorted(freq.items())
\end{lstlisting}

Ordenar por valor.

\begin{lstlisting}[language=Python]
sorted(freq.items(), key=lambda kv: kv[1])
\end{lstlisting}

Orden descendente.

\begin{lstlisting}[language=Python]
sorted(freq.items(), key=lambda kv: -kv[1])
\end{lstlisting}

\vspace{0.3cm}

\subsubsection{defaultdict}

Normalmente antes de incrementar una frecuencia se necesita comprobar si la clave existe.

\begin{lstlisting}[language=Python]
freq = {}
freq[x] = freq.get(x, 0) + 1
\end{lstlisting}

Con \texttt{defaultdict} esto desaparece.

\begin{lstlisting}[language=Python]
from collections import defaultdict

freq = defaultdict(int)

freq[x] += 1
\end{lstlisting}

También puede utilizarse para listas.

\begin{lstlisting}[language=Python]
grupos = defaultdict(list)
for x in arr:
    grupos[x % 3].append(x)
\end{lstlisting}

Con conjuntos.

\begin{lstlisting}[language=Python]
adj = defaultdict(set)
adj[u].add(v)
\end{lstlisting}

Y con cualquier función.

\begin{lstlisting}[language=Python]
from collections import defaultdict
d = defaultdict(lambda: -1)
\end{lstlisting}

\vspace{0.3cm}

\subsubsection{Counter}

Cuando únicamente interesa contar frecuencias, la clase \texttt{Counter} simplifica mucho el código.

\begin{lstlisting}[language=Python]
from collections import Counter
freq = Counter(arr)
\end{lstlisting}

Obtener el elemento más frecuente.

\begin{lstlisting}[language=Python]
freq.most_common(1)
\end{lstlisting}

Los tres más frecuentes.

\begin{lstlisting}[language=Python]
freq.most_common(3)
\end{lstlisting}

Frecuencia de un elemento.

\begin{lstlisting}[language=Python]
freq[x]
\end{lstlisting}

\vspace{0.3cm}

\subsubsection{Aplicaciones frecuentes en ICPC}

Contar frecuencias.
\begin{lstlisting}[language=Python]
freq = defaultdict(int)
for x in arr:
    freq[x] += 1
\end{lstlisting}

Agrupar elementos.
\begin{lstlisting}[language=Python]
g = defaultdict(list)
for x in arr:
    g[x % k].append(x)
\end{lstlisting}

Construir listas de adyacencia.
\begin{lstlisting}[language=Python]
adj = defaultdict(list)
for u, v in edges:
    adj[u].append(v)
\end{lstlisting}

Construir grafos sin comprobar existencia.

\begin{lstlisting}[language=Python]
adj = defaultdict(set)
adj[u].add(v)
\end{lstlisting}

Contar caracteres.

\begin{lstlisting}[language=Python]
Counter(s)
\end{lstlisting}

\vspace{0.3cm}

\subsubsection*{Resumen}

\begin{center}
\begin{tabular}{|l|l|}
\hline
Función & Uso principal \\ \hline
get() & Obtener valor sin KeyError \\ \hline
keys() & Recorrer claves \\ \hline
values() & Recorrer valores \\ \hline
items() & Recorrer clave y valor \\ \hline
update() & Actualizar con otro diccionario \\ \hline
setdefault() & Crear si no existe \\ \hline
pop() & Eliminar una clave \\ \hline
popitem() & Eliminar último elemento \\ \hline
clear() & Vaciar el diccionario \\ \hline
copy() & Copia superficial \\ \hline
len() & Número de elementos \\ \hline
defaultdict & Valor por defecto automático \\ \hline
Counter & Conteo automático de frecuencias \\ \hline
\end{tabular}
\end{center}
